% !TEX root = ../my-thesis.tex
%
\chapter{Related Work}
\label{sec:related}

This chapter reviews work that combines language models with formal or
relational argumentation and positions the present thesis against it. It first
discusses end-to-end pipelines that turn text into argumentation frameworks,
then work that uses LLMs specifically for relation extraction, the use of the
UKP corpus, and finally related formalisms.

\section{LLM-based End-to-End Argumentation Pipelines}
\label{sec:related:pipelines}

The closest work to this thesis is ARGUS~\cite{argus}, which combines LLM-based
argument mining with the construction of argumentation frameworks and symbolic
solvers in an end-to-end architecture. ARGUS shares the central idea pursued
here, namely to let an LLM handle the noisy natural-language part of the problem
and to delegate the acceptability question to a formal solver. It sets out the
architecture and the Dung-semantics perspective but does not report an empirical
evaluation of the resulting frameworks, so it leaves open how many attacks such a
pipeline accepts, whether the resulting frameworks are cyclic, and how many
extensions they admit. The present work adopts this end-to-end perspective and
supplies exactly that missing empirical layer, a reproducible implementation with
logged runs and a cross-topic analysis over eight topics.

ArgLLMs~\cite{argLLMs} also augments LLMs with formal argumentative reasoning but
takes a quantitative route. Instead of asking which arguments can be jointly
accepted, it builds Quantitative Bipolar Argumentation Frameworks (QBAFs) in
which each argument carries a numerical strength that gradual semantics propagate
along support and attack edges. The system is evaluated on binary claim
verification, where the aggregated strength of the arguments for a claim is
turned into an accept or reject decision. This differs from the present thesis in
two respects. The output is a single scalar per claim rather than a set of
jointly admissible arguments, and the evaluation targets a downstream
verification task rather than the structure of the argumentation itself.

A further recent example of an end-to-end LLM system for argument mining is the
work of Pirozelli et al.~\cite{pirozelli2026llm}, which applies LLMs to the
argument mining subtasks within a single system rather than through separate
task-specific classifiers. It reflects the broader trend that LLMs are becoming
the default engine for extraction, but like most argument mining work it stops at
the extracted structure and does not connect the result to Dung semantics or
evaluate the collective acceptability of the mined arguments. The pipeline in
this thesis differs from all three systems in its use of categorized attack
candidates with explicit justifications and confidence-based filtering, and in
computing set-based Dung extensions rather than numerical strengths or raw
extracted relations.

\section{LLMs for Attack Relation Extraction}
\label{sec:related:extraction}

The extraction step of the pipeline, which classifies the relation between every
pair of atomic arguments, is most directly related to work that studies LLMs on
relation-based argument mining. Gorur et al.~\cite{gorur2025relMining} ask
precisely whether large language models can identify the relations, such as
support and attack, that hold between arguments, the subtask on which the present
pipeline depends. Their perspective is diagnostic, in that they measure how well
LLMs recover these relations in isolation, whereas the present thesis treats
relation extraction as one stage of a larger pipeline and feeds its output into a
formal solver. Two design choices set the extraction here apart. First, the LLM
is asked not only whether an attack exists but also which of the three Pollock
types~\cite{pollock1992} it instantiates, which forces engagement with the
argumentative structure rather than a surface judgment on opposing stances.
Second, each decision carries a confidence score that routes borderline cases to
a separate verification prompt instead of accepting them outright. The value of
the typed prompt over a plain binary one is measured directly in the ablation of
Section~\ref{sec:evaluation:ablation}.

\section{Argument Mining on the UKP Corpus}
\label{sec:related:ukp}

The evaluation uses the UKP Sentential Argument Mining
Corpus~\cite{stab2018}, which provides sentence-level stance annotations for
eight controversial topics. Prior work of the supervising group applies
knowledge-based graph argument mining with topic modeling to the same
corpus~\cite{kramer2021} and reports an F1 around 67 percent. That number serves
as a reference point for the gold-standard evaluation in
Section~\ref{sec:evaluation:gold}.

\section{Related Formalisms}
\label{sec:related:formalisms}

Most closely related on the formal side is the recent work of
Alfano et al.~\cite{alfano2026}, which connects LLM-based argument mining with
argumentation and description logics through a FAKB and QBAF framework. Their
approach enriches the formal layer, both by adding logical structure through
description logics and by retaining numerical strengths through the QBAF, which
makes it more expressive but also heavier, since it commits to a strength
ordering that has to be justified. The present thesis deliberately stays with
Dung's abstract framework instead. The reason is that the question of interest
here is set-based, namely which arguments can be jointly accepted, rather than
numerical, namely how strong each argument is. Without a preference or strength
ordering, binary attack facts are the natural representation for that question,
as Chapter~\ref{sec:method} argues in more detail. The typed attacks and
confidence scores are not discarded but retained as metadata, which keeps the
door open to a richer formalism such as ASPIC+~\cite{modgil2014} or a QBAF as
future work.

\section{Positioning}
\label{sec:related:positioning}

In short, this thesis takes the end-to-end perspective of ARGUS and of the LLM
argument mining systems above, grounds it in a concrete and reproducible
implementation, and evaluates the resulting frameworks empirically, which the
closest prior systems leave undone. It differs from the gradual and
logic-augmented approaches of ArgLLMs and Alfano et al. by computing classical
Dung extensions rather than numerical strengths, and it goes beyond diagnostic
relation-extraction studies such as Gorur et al. by connecting the extracted
attacks to a formal solver. Its methodological contribution, a typed attack
extraction with confidence-based filtering, is measured directly in the ablation
of Section~\ref{sec:evaluation:ablation}.
